\documentclass[10pt,twoside]{book}
\usepackage[utf8]{inputenc}
\usepackage[T1]{fontenc}
\usepackage{geometry}
\geometry{paper=b5, margin=1.2in, top=1.5in, bottom=1.5in}
\usepackage{microtype}
\usepackage{amsmath,amssymb}
\usepackage{graphicx}
\usepackage{xcolor}
\usepackage{titlesec}
\usepackage{titletoc}
\usepackage{fancyhdr}
\usepackage{enumitem}
\usepackage{booktabs}
\usepackage{array}
\usepackage{colortbl}
\usepackage{multirow}
\usepackage[hang]{footmisc}
\usepackage{hyperref}
\hypersetup{
    colorlinks,
    linkcolor=black,
    citecolor=black,
    urlcolor=black,
}

\definecolor{chainGrey}{RGB}{80,80,80}
\definecolor{threadRed}{RGB}{160,40,40}
\definecolor{ichorBlack}{RGB}{20,20,25}
\definecolor{parchment}{RGB}{245,240,225}

\pagestyle{fancy}
\fancyhf{}
\fancyhead[RE,LO]{\textit{The Hollowed Plague -- A Sister Coven / Malachai Adventure}}
\fancyhead[LE,RO]{\thepage}
\fancyfoot[CE,CO]{\textcolor{chainGrey}{--- \textit{The witness opens its eyes. The chain tightens.} ---}}

\titleformat{\chapter}[display]
  {\normalfont\Huge\bfseries\color{threadRed}}
  {\titlerule[1pt]\vspace{4pt}\Large\chaptertitlename\ \thechapter}
  {0pt}
  {}
\titlespacing*{\chapter}{0pt}{-30pt}{20pt}

\setlength{\parindent}{0pt}
\setlength{\parskip}{0.8em}

\newenvironment{statblock}
  {\begin{center}\begin{minipage}{0.85\textwidth}\small\setlength{\parskip}{0.3em}\begin{quote}}
  {\end{quote}\end{minipage}\end{center}}

\begin{document}

\titlepage
\begin{center}
\vspace*{2cm}
{\huge\bfseries\color{threadRed} The Hollowed Plague}\\[0.8cm]
{\Large\itshape A War Between the Sister Coven and the Chained Angel}\\[1.5cm]
\rule{0.6\textwidth}{1pt}\\[0.5cm]
\emph{An Epic Adventure for Tier III--IV Characters}\\[1cm]
\vfill
{\small \textcolor{chainGrey}{Under the lamp, bread and salt --- and a witness that no longer sleeps.}}\\[0.2cm]
{\small \textcolor{chainGrey}{The ink is fresh. The road forks. The Hollowed watches.}}
\end{center}
\newpage

\chapter*{Safety and Consent at the Table}
\addcontentsline{toc}{chapter}{Safety and Consent at the Table}
This adventure includes themes of body horror (hollowed creatures, undeath), loss of identity, coercion, and moral choices with no clean solution.

Before play, establish:
\begin{itemize}
\item \textbf{Lines:} Hard boundaries. Example: no sexual violence, no harm to children (the Hollowed are child‑like but are not human children), no graphic torture.
\item \textbf{Veils:} Soft boundaries. Example: the transformation of a Hollowed can be described obliquely rather than depicted in detail.
\item \textbf{X‑Card:} Any player may tap an X card to pause or remove a scene without explanation. The table pivots immediately.
\end{itemize}
Use these tools. Trust them.

\tableofcontents
\newpage

\chapter{Adventure Overview}

\section{The Hollowed Plague}
The Daughters of the Unbroken Cord have maintained the Hollowed for centuries -- heartless children with silver wire eyes, perfect witnesses to every deathbed oath and whispered debt. They do not speak. They do not hunger. They simply watch.

Now something has changed.

Malachai, the Chained Angel, has found a way to curse the Hollowed. The curse does not destroy them; it \textit{wakens} them. They gain free will. They remember the debts they witnessed -- and they resent every one. The first Hollowed to turn did so in the hospice of Sister Aris. It rose from its corner, blinked its silver eyes, and walked out into the world.

Now they stalk the countryside. They are not mindless. They are witnesses turned persecutors. They seek out those who owe debts to the Chain -- and they collect in blood.

The Daughters are paralysed. Their own witnesses have become their hunters. The High Mother weaves furiously, trying to contain the corruption, but the Cord frays with each Hollowed that breaks free.

\section{The Fork -- Two Paths, Two Cures, One Choice}
At the end of Act I, the party must choose how to stop the plague. There is no third way -- only the consequences of the path taken.

\begin{itemize}
\item \textbf{The House Without Doors (Chain Path):} Retrieve a \textbf{First Needle}, one of the original artifacts used to hollow the first child. Using the Needle will \textbf{end the plague} and return the Turned Hollowed to their original state -- silent, heartless witnesses. They will not be free, but they will not hunt. The Chain survives, but the Hollowed remain prisoners.
\item \textbf{Zahk"Kozahn (Malachai Path):} Steal a \textbf{Flash of Ichor}, a drop of Malachai's own blood. Using the Ichor will \textbf{give the Turned Hollowed true sentience} -- they will be free, with wills of their own. But Malachai's gift is never clean. The Ichor carries a hidden curse. The Hollowed may be alive, but they will always hunger for something. No one knows what.
\end{itemize}

\section{The Hidden Consequence -- Doing Nothing}
If the party delays too long or refuses to choose, the war escalates. The Turned Hollowed, desperate and angry, seek out the source of their suffering -- the High Mother. They assault the House Without Doors, forcing her to unravel the Cord in self‑defence. The resulting explosion of grey thread does not stop the Hollowed; it scatters them across every threshold, every doorway, every birth and death bed in the known world. The Hollowed become a permanent, haunting plague -- witnesses without masters, cursed to remember every sin they ever saw.

\textit{The Cord breaks. The nations burn. There is no third option.}

\section{The Corruption Clock -- Hollowed Awakening}
A single 8‑segment clock tracks the spread of the curse among the Hollowed.

\textbf{Hollowed Awakening [8]:} Ticks when:
\begin{itemize}
\item A Hollowed attack is witnessed (or a corpse found). +1
\item A Daughter is killed by a Hollowed. +2
\item The party fails to protect a Daughter or a hospice. +1
\item A PC willingly accepts a curse from Malachai. +2
\end{itemize}
At 4 segments: Hollowed begin to form packs; travel at night is dangerous.  
At 6 segments: A Sister coven is overrun; a Winter is killed.  
At 8 segments: The High Mother begins to unravel; the House Without Doors begins to collapse. The party must resolve the adventure before the clock fills or face the cataclysmic outcome (the Cord breaks, nations burn).

\section{Recommended Tier}
III--IV (experienced characters with assets, followers, and Patron ties).

\section{Estimated Play Time}
3--4 sessions.

\chapter{Act I -- The Awakening}

\section{The Hook}
The party is travelling through a rural valley when they find a small hospice -- a Daughter's waystation -- abandoned, its doors torn open. Inside, the bodies of two Daughters lie twisted, their chests hollowed out. Their eyes are missing.

A third Daughter, gravely wounded, clutches a grey cord. "The Hollowed," she whispers. "They... woke up. They asked us why we took their hearts. Then they took ours."

She dies.

In her hand is a silver coin -- warm, stamped with a ship. A Malachai coin.

\section{Key NPCs -- Act I}
\begin{itemize}
\item \textbf{Sister Vorna (Winter of the Chain):} A severe woman with grey braids and a Thorn‑Pin in her hair. She is the only Winter in the region not yet attacked. She suspects Malachai is behind the curse, but she cannot prove it. She needs the party to investigate. \textit{"The Hollowed are our eyes. If they have turned, we are blind. Find me the source -- the first needle, the first coin -- before we are all deaf as well."}
\item \textbf{The Hollowed (First Breaker):} A child with silver wire eyes, but its mouth is no longer sewn shut. It speaks in a flat, echoing voice. \textit{"The Chain stole my heart. The Angel gave me a voice. Why should I watch for you when I never chose to see?"} This Hollowed can be negotiated with, fought, or studied. It carries a shard of ichor (not enough to cure, but enough to prove Malachai's involvement).
\item \textbf{"Broken" Lysias (Malachai Cultist, optional):} A Malachai follower who has fallen out of favor. He offers to trade information on the curse for protection. He is terrified of what he helped unleash.
\end{itemize}

\section{Key Encounters -- Act I}

\subsection{The Waking Hollowed (Investigation / Social)}
The party tracks the First Breaker to an abandoned mill. The Hollowed is not hostile -- it is confused, angry, and desperate for answers. It asks:
\begin{itemize}
\item "Who made me?"
\item "Why am I not allowed to speak?"
\item "What is the debt I witnessed? No one told me why I had to watch."
\end{itemize}
The party can learn that the curse was planted when a Malachai cultist placed a silver coin in the mouth of a Hollowed (a secret ritual). Success (Wits + Insight, DV 4) reveals that the coin was not just any coin -- it was a \textbf{First Coin}, blessed by Malachai herself.

\subsection{The First Coin (Recovery)}
The coin is still in the mill, lodged in the Hollowed's chest. To retrieve it, the party must:
\begin{itemize}
\item Negotiate: Presence + Sway, DV 4. The Hollowed wants to know what it will gain by giving up its "gift."
\item Combat: The Hollowed (Cap 3, Tier III) attacks. It is fast, strong, and can phase through shadows. Defeating it yields the coin.
\item Stealth: While another PC distracts it, a thief can extract the coin (Wits + Subterfuge, DV 5).
\end{itemize}
The coin is warm and pulses. If a PC spends it (gains 2 Boons, but begins a Curse Clock), they become a target for the Hollowed -- and also gain a clue about the location of the First Needle or the Ichor.

\subsection{The Winter's Judgment}
Sister Vorna examines the coin. "This is not a normal curse. It is an anchor. Malachai has attached something to the Hollowed -- a thread that leads back to... somewhere. I cannot see the end. But you can."

She offers two paths. "The First Needles are kept in the House Without Doors, in the Chamber of the First Hollowed. No Daughter may enter -- the Runekeepers of the High Mother guard it as their highest charge. But you are not Daughters. You may pass if you can survive the liminal threshold."

"Or you could trace the coin to its source -- to the cursed city, to Zahk"Kozahn. But you cannot enter that place without bearing a curse yourself. You would have to accept a Malachai gift. That is a price I cannot ask you to pay."

"Choose quickly. The Hollowed are waking faster. If you do nothing... the Cord will break. And when the Cord breaks, every threshold will bleed. Nations will burn."

\section{Act I Resolution}
The party leaves the hospice with the coin and the knowledge of the two paths.
\begin{itemize}
\item \textbf{Advance the Hollowed Awakening clock} by +1 for each Daughter killed in the hospice.
\item If they killed the First Breaker, advance by +1 (the violence attracts more Hollowed).
\end{itemize}
The GM then presents the fork. The party must choose. There is no turning back.

\chapter{Act II -- The Fork}

\section{Path A -- The House Without Doors}

\subsection{Requirement}
No special curse needed. The party must reach a liminal threshold -- a crossroads at dawn, a river that flows uphill, or a door that appears only when a specific debt is spoken aloud. Sister Vorna provides the ritual to open the way.

\subsection{Step 1 -- The Liminal Trial (Threshold Challenges)}
The party must pass three challenges, each representing a trial of the Chain's ethos. Failure advances a \textbf{Threshold Collapse [6]} clock; when full, the way is blocked and the party must find another entrance (costs time, ticks the Hollowed Awakening clock by +2).

\subsubsection*{Trial of the Unspoken Name}
A door asks: "What is the name of the first child who was hollowed?" The answer is not written anywhere. It is a secret held only by the High Mother. The party must either:
\begin{itemize}
\item \textbf{Research:} Wits + Lore, DV 5, using fragments of lore from the hospice (interviews with Daughters, old ledgers). Success gives a name: \textit{Ellinor, the Stillborn Daughter}.
\item \textbf{Bargain:} A Silent One (a Hollowed not yet turned) appears and offers the name in exchange for a memory (permanent, choose a minor fact).
\item \textbf{Guess:} Wrong answers trigger a \textbf{Wrath of the Unnamed} -- the door attacks with psychic harm (Resist DV 3, 2 Fatigue on failure).
\end{itemize}

\subsubsection*{Trial of the Witness}
The party crosses a hall lined with mirrors. In each reflection, they see a version of themselves who accepted a Malachai coin. The mirrors ask: "What would you give up to be free of that temptation?" Each PC must answer. Honest answers (roleplayed, no roll) let them pass. Deflection or lies cause a mirror to shatter, releasing a \textbf{Reflection Wraith} (Cap 2, Tier II) that attacks.

\subsubsection*{Trial of the Unbroken Cord}
A thread of grey cord spans a chasm. To cross, the party must tie themselves together with a piece of cord (they have to use their own bonds or a symbolic thread). The cord tests their loyalty: if any PC acts selfishly (tries to cross alone, hoards a resource), the cord snaps and the party is separated (lose 1 Segment on the Threshold Collapse clock). If they cooperate, they gain a \textbf{Cord Token} (once, re‑roll a failed social roll with a Daughter).

\subsection{Step 2 -- The Siege of the House}
The House Without Doors is not a building -- it is a shifting, liminal space. As the party approaches, they find it besieged by \textbf{Turned Hollowed} -- infected witnesses, their silver eyes now black, their mouths open in silent screams. The Hollowed are trying to break in, not to destroy the House, but to \textit{enter} -- to confront the High Mother.

The party must fight or sneak through the siege to reach the threshold.

\textbf{Siege Clock [6]:} Ticks as the party fights. At 2, reinforcements arrive (more Turned Hollowed). At 4, a \textbf{Hollowed Pack Leader} (Cap 4, Tier III) appears. At 6, the party is overwhelmed -- they must flee or find a different entrance (the House may open a side door for a price).

\textbf{Encounter Types:}
\begin{itemize}
\item \textbf{Stealth:} Wits + Stealth, DV 4, to slip between patrolling Hollowed. Failure triggers combat.
\item \textbf{Combat:} 2--4 Turned Hollowed (Cap 2 each, Tier II). They have silver wire eyes and clawed hands. They are silent.
\item \textbf{Distraction:} Presence + Deception, DV 3, to lure a group away. Success reduces the Siege Clock by 1.
\end{itemize}

\subsection{Step 3 -- The Guardians of the Needle (Final Encounter)}
Inside the House, the party finds the \textbf{Chamber of the First Hollowed}. A dozen \textbf{Runekeeper Daughters of the High Mother} guard it -- the most loyal and powerful members of the Chain, their flesh codices glowing.

They do not attack immediately. The eldest speaks: "You come for the First Needle. We know. The High Mother knows. She will not give it freely. To take it, you must prove you are not here to serve the Chained Angel."

\textbf{The Trial of the Needle:}
The Runekeepers demand the party answer three riddles (related to debt, witness, and mercy). Each riddle is a test of the Chain's philosophy. The party may:
\begin{itemize}
\item \textbf{Answer correctly:} Presence + Lore, DV 4 per riddle. Success grants a boon; failure forces a fight.
\item \textbf{Refuse and fight:} Combat against 2--3 Runekeepers (Cap 4 each, Tier III--IV). They use \textit{Cord of Flesh} (store Fatigue) and \textit{Unnamed Strike} (temporary name loss).
\item \textbf{Bargain:} Offer a memory, a name, or a future debt to the High Mother in exchange for the Needle. The Runekeepers accept if the offer is weighty enough (e.g., a PC's true name, a bond with a loved one).
\end{itemize}

Once the Needle is retrieved (a slender, cold silver rod, warm to the touch), the House begins to shake. The High Mother weeps. "Go. HEAL what I broke." A door opens to the outside -- directly into the valley where the adventure began.

\subsection{Path A Resolution}
The party escapes with the \textbf{First Needle}. They can now attempt to end the plague by returning the Hollowed to their original state (see Act III, Resolution A).

\section{Path B -- Zahk"Kozahn}

\subsection{Requirement}
At least one PC must be cursed with a Malachai gift (vampirism, lycanthropy, a thorn, a silver coin used). If none are, they must accept a curse from a Malachai cultist (or from the First Breaker) -- this begins a Curse Clock.

\subsection{Step 1 -- The Road to the Cursed City}
Zahk"Kozahn cannot be found by the uncursed. A cursed PC leads the party through a dream‑scape of drifting sands and broken fig trees. The journey takes three days, but in the waking world only one night passes.

\textbf{Travel Clocks:}
\begin{itemize}
\item \textbf{Desert Hunger [6]:} Ticks as the party suffers the curse's influence. At 2, the cursed PC must feed (take 1 Fatigue or attack an ally). At 4, hallucinations begin (--1 die to all rolls). At 6, the cursed PC transforms partially (gain a mutation).
\item \textbf{Approach to the City [8]:} Ticks up as the party progresses. Must fill to reach the gates.
\end{itemize}

\subsection{Step 2 -- The Ghoul City}
Zahk"Kozahn is not deserted. It is populated by \textbf{Worn} -- the ghoul descendants of those who served Malachai before her chaining. They are undead, desperate, and strung out on ichor residue.

The party must cross several districts to reach the ziggurat.

\textbf{Key Encounters:}
\begin{itemize}
\item \textbf{The Market of Whispers:} Ghouls trade memories for ichor. The party can barter (give a memory, gain a clue) or fight (the ghouls are Cap 2 each, Tier II, but numerous). Success reveals the path to the ziggurat.
\item \textbf{The Fig Tree Grove:} A grove where cursed trees bear fruit that whispers secrets. Eating a fruit gives a vision of the Well of Ichor (grants +1 die to the final descent) but marks 1 Corruption.
\item \textbf{The Silent Procession:} A parade of Hollowed (Turned) marches through the city, led by a \textbf{Collector} (a child with silver wire eyes, now grown). They do not attack unless provoked. If the party hides (Wits + Stealth, DV 5), they pass safely. If they fight, they face a difficult battle (Cap 5, Tier IV).
\end{itemize}

\subsection{Step 3 -- The Well of Ichor (Final Encounter)}
Beneath the ziggurat, in the cathedral where Malachai is chained, the party finds the \textbf{Well}. The ichor is black, thick, and steaming. Malachai's voice whispers from the chains: "Take what you need. But know: a flash of ichor is a fragment of my blood. It will remember you. And so will they."

To extract a \textbf{Flash of Ichor} (a single drop, sealed in a phial), the party must:
\begin{itemize}
\item \textbf{Drink from the Well (cursed PC only):} They gain 2 Boons and a permanent Ichor‑Touched boon (once per session, ignore a curse effect), but the Flash is collected automatically.
\item \textbf{Use a ritual:} Arcana + Lore, DV 5. Success extracts a drop without drinking. Failure causes the ichor to spatter; the cursed PC must test Resolve (DV 4) or gain a temporary mutation (e.g., silver eyes, thirst for blood).
\item \textbf{Steal from Malachai's chains:} A direct confrontation with the Chained Angel (she cannot move, but she can curse). Each PC must test Spirit + Resolve vs DV 5 or gain a Curse Clock segment. Success allows one PC to chip a fragment of dried ichor from her binding.
\end{itemize}

Once the Flash is secured, Malachai laughs. "You think you can give them freedom? You can. But freedom is a hunger. I would know."

A door opens -- a threshold back to the waking world. The party must escape before the curse of the city claims them (run a chase; failures tick the Desert Hunger clock).

\subsection{Path B Resolution}
The party escapes with the \textbf{Flash of Ichor}. They can now attempt to give the Turned Hollowed true sentience (see Act III, Resolution B).

\chapter{Act III -- The Debt Called Due}

The party has either the First Needle or the Flash of Ichor (or, perhaps miraculously, both). They must return to a Daughter's hospice -- or any threshold sanctified by the Chain -- and perform a ritual.

\section{The Ritual}
Sister Vorna (or a surviving Winter) prepares the space. The party must protect her while she works.

\textbf{Ritual Clock [8]:} Ticks as the ritual progresses. Each exchange, a PC can contribute (choose a skill: Arcana to guide the energy, Presence to soothe the Cord, Body to hold the threshold steady). On a success, tick the clock; on a failure, the GM ticks the \textbf{Hollowed Awakening} clock instead (the curse spreads faster).

\subsection{Complications During the Ritual}
At each segment tick, roll 1d6 on the table:
\begin{enumerate}
\item A Turned Hollowed attacks the hospice. (Combat, Cap 3)
\item The ritual backfires -- one PC gains a Curse Clock segment.
\item A Silent One (still loyal) appears and offers to help, but demands a memory.
\item The Cord frays -- a Daughter falls to the curse; she becomes a Turned Hollowed.
\item Malachai's voice echoes; all cursed PCs test Resolve DV 4 or lose their next action.
\item A flash of ichor spills; the ground becomes difficult terrain (--1 Position to all).
\end{enumerate}

\section{Resolution A -- Using the First Needle (End the Plague, Return Hollowed to Their Original State)}
The needle is inserted into the Cord's heart (a glowing knot at the centre of the ritual circle). The High Mother's scream shakes the House. All Turned Hollowed in the region collapse, their eyes closing. They are no longer cursed -- they are simply dead. The Hollowed that remain are restored to their original, silent, heartless state. They will witness again. They will never wake.

The Chain survives. The Daughters weep with relief.

\textbf{Consequences:}
\begin{itemize}
\item The Hollowed Awakening clock stops.
\item The party gains a permanent String: \textbf{Favor of the High Mother} -- once per session, reroll a failed roll involving memory or oath.
\item Malachai is infuriated; she sends a Collector after the party (future adventure hook).
\item The Hollowed are not free. They will never be free. This is mercy to the Chain, but is it mercy to them?
\end{itemize}

\section{Resolution B -- Using the Flash of Ichor (Give the Turned Hollowed Sentience)}
The ichor is dripped onto the Cord. It hisses and smokes, and for a moment, the Cord turns black. Then it pulses -- and the curse does not reverse. It \textit{transforms}. The Turned Hollowed do not return to their silent selves. They become fully awake -- conscious, aware, and free. Their silver wire eyes now gleam with intelligence. They speak as one voice: "We see. We remember. We choose."

But Malachai's gift is never clean. The ichor carries a hidden curse. The Hollowed are free, but they now \textit{hunger} -- not for blood, but for truth. They will seek out every debt they ever witnessed and demand restitution. Some will be merciful. Others will not.

The Hollowed are no longer prisoners. They are a new people -- born of suffering, shaped by debt, and armed with every secret the dying ever whispered.

\textbf{Consequences:}
\begin{itemize}
\item The Hollowed Awakening clock stops.
\item The party gains a permanent String: \textbf{Ichor Debt} -- once per session, call on a Hollowed for a single piece of witnessed information. But the Hollowed may also call on you.
\item The Daughters are horrified. The High Mother weaves frantically, trying to re‑establish control. She will fail.
\item The party has given freedom to the enslaved. But freedom comes at a price, and the invoice is still being written.
\end{itemize}

\section{Resolution C -- Doing Nothing (The Cord Breaks, Nations Burn)}
If the party refuses to choose, or if the Hollowed Awakening clock fills before the ritual is attempted, the House Without Doors collapses. The High Mother, in her desperation, unravels the Cord to prevent the Hollowed from entering. The resulting explosion of grey thread does not stop the Hollowed -- it scatters them across every threshold, every doorway, every birth and death bed in the known world.

The Hollowed become a permanent, haunting plague -- witnesses without masters, cursed to remember every sin they ever saw. They appear at funerals, at weddings, at the bedsides of the dying. They do not speak. They only watch. And everyone knows that they are judging.

\textit{The Cord breaks. The nations burn. There is no third option.}

\textbf{Consequences:}
\begin{itemize}
\item The campaign shifts: the party must now find a way to either destroy the Hollowed or broker a peace between them and the living.
\item Every settlement now has a Hollowed problem. Travel is dangerous. Trust is eroded.
\item Malachai whispers: "You could have chosen. You chose nothing. Now everyone pays."
\end{itemize}

\chapter{GM Tools}

\section{Key NPC Stat Blocks (Simplified)}
\subsection{Turned Hollowed (Cap 2--3, Tier II--III)}
\textit{Abilities:} Silver Gaze (target must test Resolve or reveal a secret), Silent Step (+2 dice to stealth), Claw (+1 Harm). \textit{Weakness:} Can be temporarily pacified by speaking a name they once witnessed (requires knowledge of the memory).

\subsection{Runekeeper Daughter of the High Mother (Cap 4, Tier III--IV)}
\textit{Abilities:} Cord of Flesh (store Fatigue as Knots), Unnamed Strike (target loses name temporarily), Cord Pull (teleport ally). \textit{Weakness:} Cannot cross a threshold of running water.

\subsection{Malachai's Ghoul (Worn) (Cap 2, Tier II)}
\textit{Abilities:} Ichor Hunger (feed on a corpse to heal 1 Harm), Desperate Lunge (extra action when below half health). \textit{Weakness:} Fire (takes +1 Harm).

\section{Clocks Summary}
\begin{itemize}
\item \textbf{Hollowed Awakening [8]} -- Global curse pressure.
\item \textbf{Threshold Collapse [6]} (House path) -- time to enter.
\item \textbf{Siege Clock [6]} (House path) -- battle outside the House.
\item \textbf{Desert Hunger [6]} (Zahk"Kozahn path) -- curse attrition.
\item \textbf{Approach to the City [8]} (Zahk"Kozahn path) -- travel to the city.
\item \textbf{Ritual Clock [8]} (final) -- success or failure.
\end{itemize}

\section{Random Complication Table (d12)}
\begin{enumerate}
\item A Hollowed breaks free of the siege and attacks a PC carrying the Needle.
\item The desert hunger manifests: a cursed PC must feed or take Fatigue.
\item A Runekeeper Daughter recognises a PC from a past debt.
\item Malachai whispers a tempting offer to a PC (gain 2 Boons, start a Curse Clock).
\item The Cord frays; a Daughter dies, and her Hollowed turns.
\item A Silent One (not yet turned) appears and offers critical information for a memory.
\item The ritual site is a former graveyard; the dead rise to interfere.
\item A Collector demands a name as payment for safe passage.
\item The fig tree fruit is hallucinogenic; a PC loses track of time.
\item A ghoul marketeer recognises the ichor phial and tries to steal it.
\item The High Mother speaks directly to the party: "Choose quickly. I am fraying."
\item A second Flash of Ichor is found -- but it costs a finger to retrieve.
\end{enumerate}

\chapter{Colophon}
\addcontentsline{toc}{chapter}{Colophon}
This adventure was written in a hospice where a Hollowed had just turned. The Sister there did not scream. She simply tied a grey cord around her own wrist and waited.

She is still waiting.

\vspace{0.5cm}
\noindent
\textit{The ink is fresh. The road forks. The witness opens its eyes. Choose.}

\end{document}
